%%%%%
%
\section{Temp�rature du point de ros�e � partir de
l'humidit� absolue}

\itemstyle{Op\'erateur~:} \operfmt{Td\_Ha}
\index{Td\_Ha}

%%%%%OBJET%%%%
\itemstyle{Objet~:}
Calcul de la temp�rature du point de ros�e � partir de
l'humidit� absolue et de la temp�rature de l'air.

%%%%%ARGUMENTS%%%%
\bigskip
%%%%%ENTREES%%%%
\begin{tabarg}
\multicolumn{3}{l}{\titretab {Entr\'ees}}\\
\hline
$Ha$ & Vecteur, \texttt{float} & Humidit� absolue (g/m$^3$)\\
$T$ & Vecteur, \texttt{float} & Temp�rature de l'air (Celsius ou K)\\
\end{tabarg}

\medskip
%%%%%SORTIES%%%%
\begin{tabarg}
\multicolumn{3}{l}{\titretab {Sorties}}\\
\hline
$Td$ & Vecteur, \texttt{float} & Temp�rature du point de ros�e (Celsius)\\
\end{tabarg}


%%%%%FORMULE%%%%
\itemstyle{Formule}

%%Td_Ha
T est converti en Kelvin, si besoin est, puis :
\[
e = \frac{Ha}{1000} . Rv . T \qquad \text{(avec } Rv  = 461.51  J K^{-1} kg^{-1} \text{)}
\]
puis 
\newcommand{\InvEwBoltonRes}{\ensuremath{Td_K}}
%% Formule (et uniquement la formule) de la fonction
%% inv_ew_bolton() de la biblioth�que thermo.

%La commande \providecommand permet de ne prendre en
%compte la nouvelle commande que si elle n'existe pas
%d�j�.

\providecommand{\InvEwBoltonRes}{T}
\[
\InvEwBoltonRes = \frac{29.65 * \ln(e/(100*6.112)) - 4826.56}{\ln(e/(100*6.112)) - 17.67}
\]


et $Td  = \InvEwBoltonRes -273,15 $ (conversion en degr�s Celsius).


%%%%%COMMENTAIRE%%%%
\itemstyle{Commentaire}

$ e = ew(Td)$ (d�finition de Td) et $Ha = e / (Rv . T)$ (Dalton et Gaz Parfait) 
s'inversent en  $Td = ew^{-1}(Ha . Rv . T)$.

%%%%%AUTEUR%%%%
\itemstyle{Auteur~:}
CNRM/GMEI/TRAMM

%%%%%REFERENCE%%%%
\itemstyle{R\'ef\'erences~:}
La fonction r�ciproque de $e_w(T)$ utilis�e est celle de Bolton-1980 \cite{Bolton-80}
